% Part: first-order-logic
% Chapter: syntax-and-semantics
% Section: introduction

\documentclass[../../../include/open-logic-section]{subfiles}

\begin{document}

\olfileid{fol}{syn}{int}

\olsection{ভূমিকা}

প্রথম-ক্রম যুক্তিবিদ্যার তত্ত্ব ও অধিতত্ত্ব গড়ে তুলতে হলে প্রথমেই তার
রাশিগুলির সংকেতবিন্যাস ও অর্থতত্ত্ব সংজ্ঞায়িত করতে হবে। প্রথম-ক্রম
যুক্তিবিদ্যার রাশি হলো পদ ও !!{formula}s। !!{variable}s,
!!{constant}s ও !!{function}s থেকে পদ গঠিত হয়। আবার !!^{formula}s,
!!{predicate}s-এর সঙ্গে পদ মিলিয়ে সবচেয়ে ছোট, ``পরমাণু''
!!{formula}s গঠিত হয়; তারপর যৌক্তিক সংযোজক ও পরিমাণসূচক ব্যবহার করে
পরমাণু !!{formula}s থেকে আরও জটিল সূত্র গঠন করা যায়। গঠনের নিয়মগুলি
নির্ধারণের নানা উপায় আছে; আমরা তার মধ্যে কেবল একটি উপায় দিচ্ছি।
অন্য ব্যবস্থায় ভিন্ন সংকেত বেছে নেওয়া হতে পারে, ভিন্ন সংযোজকগুচ্ছকে
মৌলিক ধরা হতে পারে, বন্ধনী ভিন্নভাবে ব্যবহার করা হতে পারে (কিংবা
তথাকথিত পোলিশ সংকেতরীতির মতো একেবারেই ব্যবহার না-ও করা হতে পারে)।
তবে সব পদ্ধতির সাধারণ বৈশিষ্ট্য হলো, গঠনের নিয়মগুলি পদ ও
!!{formula}s-এর সেটকে \emph{আরোহীভাবে} সংজ্ঞায়িত করে। নিয়মগুলি
যথাযথ হলে, প্রত্যেকটি রাশি মূলত একটিমাত্র উপায়েই সেই নিয়ম মেনে
গঠিত হতে পারে। আরোহী সংজ্ঞা থেকে রাশিগুলি \emph{একক পাঠযোগ্য}
হওয়ার ফলে একই পদ্ধতিতে—আরোহী সংজ্ঞার সাহায্যে—আমরা সেগুলির অর্থও
দিতে পারি।

রাশির অর্থ দেওয়া অর্থতত্ত্বের বিষয়। অর্থতত্ত্বের কেন্দ্রীয় ধারণা
হলো !!a{structure}-এ পরিতৃপ্তি। !!^a{structure} ভাষার
নির্মাণ-উপাদানগুলিকে অর্থ দেয়: !!a{domain} হলো বস্তুর একটি অশূন্য
সেট। পরিমাণসূচকগুলিকে এই সংজ্ঞাক্ষেত্রের উপর বিস্তৃত বলে ব্যাখ্যা
করা হয়, !!{constant}s-কে সংজ্ঞাক্ষেত্রের সদস্য আরোপ করা হয়,
!!{function}s-কে তাদের স্থানসংখ্যা অনুযায়ী
!!{domain}-এর সসীম টিউপল থেকে সংজ্ঞাক্ষেত্রে ফলন আরোপ করা হয়,
এবং !!{predicate}s-কে !!{domain}-এর উপর সম্বন্ধ আরোপ করা হয়।
% BN-SRC-501 TARGET-CORRECTION-BEGIN
\par\smallskip\noindent\textnormal{\small\textbf{উৎস-সংশোধন (BN-SRC-501):} মূল সারাংশে সব অপেক্ষকচিহ্নকে সংজ্ঞাক্ষেত্র থেকে নিজের মধ্যেই এক-স্থানের অপেক্ষক দেওয়া হয় বলে মনে হয়। এই অধ্যায়ের আনুষ্ঠানিক গঠন-সংজ্ঞা $n$-স্থানের চিহ্নে সংজ্ঞাক্ষেত্রের $n$-টিউপল থেকে সংজ্ঞাক্ষেত্রে অপেক্ষক দেয়; বাংলা সারাংশে সেই প্রযোজ্য স্থানের সংখ্যা স্পষ্ট করা হয়েছে।}
% BN-SRC-501 TARGET-CORRECTION-END
!!{domain} এবং মৌলিক শব্দভাণ্ডারে করা আরোপগুলি একত্রে
!!a{structure} গঠন করে। !!^{variable}s, !!{formula}s-এ থাকতে পারে;
তাদের অর্থতত্ত্ব দিতে !!{domain}-এর !!{element}s-ও তাদের আরোপ করতে
হয়—এটিই চলরাশি-আরোপ। শেষে পরিতৃপ্তি-সম্বন্ধ এই সবকিছুকে একত্র করে।
কোনো !!^a{formula}, !!a{structure}~$\Struct{M}$-এ চলরাশি-আরোপ~$s$-এর
সাপেক্ষে পরিতৃপ্ত হতে পারে; একে $\Sat{M}{!A}[s]$ লিখি। এই
সম্বন্ধটিকেও~$!A$-এর গঠনের উপর আরোহ করে সংজ্ঞায়িত করা হয়। যেমন,
$!A \land !B$-এর পরিতৃপ্তিকে $!A$ ও~$!B$-এর পরিতৃপ্তি (বা
অপরিতৃপ্তি) দিয়ে সংজ্ঞায়িত করতে যৌক্তিক সংযোজকগুলির সত্যক-সারণি
ব্যবহার করি। তারপর দেখা যায়, !!{formula}~$!A$ যদি !!a{sentence}
হয়, অর্থাৎ তার কোনো মুক্ত চলরাশি না থাকে, তাহলে চলরাশি-আরোপটি
অপ্রাসঙ্গিক। ফলে !!{sentence}s, !!{structure}s-এ কেবল পরিতৃপ্ত (বা
অপরিতৃপ্ত)—এ কথাই বলতে পারি।

বাক্যের জন্য পরিতৃপ্তি-সম্বন্ধ $\Sat{M}{!A}$-এর ভিত্তিতে
বৈধতা, ফলশ্রুতি ও পরিতৃপ্তিযোগ্যতার মৌলিক অর্থতাত্ত্বিক ধারণাগুলি
সংজ্ঞায়িত করতে পারি। কোনো বাক্য বৈধ, $\Entails !A$, যদি প্রত্যেক
গঠন তাকে পরিতৃপ্ত করে। !!{sentence}s-এর কোনো সেট থেকে তার ফলশ্রুতি
হয়, $\Gamma \Entails !A$, যদি $\Gamma$-র সব !!{sentence}s-কে
পরিতৃপ্ত করে এমন প্রত্যেক !!{structure}, $!A$-কেও পরিতৃপ্ত করে।
আর বাক্যের কোনো সেট পরিতৃপ্তিযোগ্য, যদি কোনো !!{structure} সেটটির সব
!!{sentence}s-কে একই সঙ্গে পরিতৃপ্ত করে। !!{formula}s আরোহীভাবে
সংজ্ঞায়িত এবং পরিতৃপ্তিও পালাক্রমে !!{formula}s-এর গঠনের উপর আরোহ
করে সংজ্ঞায়িত বলে আমাদের অর্থতত্ত্বের ধর্মগুলি প্রমাণ করতে এবং
সংজ্ঞায়িত অর্থতাত্ত্বিক ধারণাগুলির মধ্যে সম্বন্ধ স্থাপন করতে আমরা
আরোহ ব্যবহার করতে পারি।

\end{document}
